V \myref{pododdíle}{sssec:relace} jsme definovali
\hyperref[def:relace-ekvivalence]{relaci ekvivalence}. Je to relace, která v
podstatě umožňuje považovat dva prvky množiny za \emph{stejné}, za předpokladu,
že sdílejí jisté vlastnosti. Intuitivně využíváme ekvivalence neustále, téměř
kdykoli pojmenováváme objekty. Žádné dva objekty, které vnímáme, nejsou přímo
\emph{stejné}, ale přiříkáme jim tatáž jména, protože jsou \uv{dostatečně
podobné}, abychom je mohli označit stejně. Kupříkladu, starou dřevěnou vysokou
židli i novou malou kovovou židli jsme právě nazvali \emph{židlí}, neboť obě
mají dostatek společných charakteristik a plní stejný účel, přestože se jistě
nejedná o týž předmět.

Příklad židle s sebou nese další samozřejmou, leč užitečnou, myšlenku: svůj svět
přirozeně rozdělujeme na skupiny předmětů, které považujeme za \uv{stejné}. V
matematice se taková \emph{skupina podobných předmětů} nazývá \alert{třída
ekvivalence}. Je to množina všech prvků, které jsou vzájemně ekvivalentní.

\begin{definition}[Třída ekvivalence]
  \label{def:trida-ekvivalence}
  Ať $e$ je ekvivalence na množině $a$. Pro libovolný $x \in a$ říkáme všem
  prvkům z $a$, které jsou ekvivalentní s $x$, \emph{třída ekvivalence příslušná
  $x$} a značíme ji $[x]_e$. Formálně, pro každé $x \in a$ je
  \[
    [x]_e = \{y \in a \mid (x, y) \in e\}.
  \]
  Prvku $x$ pak říkáme \emph{representant} této třídy ekvivalence.
\end{definition}
\begin{example}
  \label{exam:trida-ekvivalence-1}
  Ať $l$ je množina lidí a $e$ je ekvivalence \uv{být stejně starý} (je to
  opravdu ekvivalence?). Když je $\text{Adam} \in l$ člověk, kterému je 29 let.
\end{example}
